%% ---------------------------------------------------------------
%%  Pager – Proyecto Final de Ciencias de la Computación – UTEC
%%  Segovia Ureta, Gianpier | Lazo Mera, Kalos
%%  Asesor: Martinez Abaunza, Victor Eduardo
%% ---------------------------------------------------------------
\documentclass[a4paper,11pt]{article}

%% Paquetes mínimos
\usepackage[utf8]{inputenc}
\usepackage[T1]{fontenc}
\usepackage[spanish]{babel}
\usepackage[top=2.2cm,bottom=2.2cm,left=2.5cm,right=2.5cm]{geometry}
\usepackage{hyperref}
\usepackage{booktabs}      % tablas más elegantes
\usepackage{array}
\usepackage{parskip}       % espaciado entre párrafos en lugar de sangría
\usepackage{titlesec}      % control de secciones
\usepackage{enumitem}
%% Reducir espacio vertical de secciones
\titlespacing*{\section}{0pt}{8pt}{4pt}
\titlespacing*{\subsection}{0pt}{6pt}{3pt}

\hypersetup{colorlinks=true, urlcolor=blue, linkcolor=black}

%% ---------------------------------------------------------------
\begin{document}

%% Encabezado manual (sin maketitle para ahorrar espacio)
\begin{center}
  {\large\bfseries Mitigación de la Resistencia Evolutiva en Glioblastoma\\
  Multiforme mediante MARL Adversarial: Optimización de\\
  Políticas Adaptativas bajo un Framework MAPPO-CTDE\\
  y Simulación Calibrada}\\[8pt]
  {\small Segovia Ureta, Gianpier \quad|\quad Lazo Mera, Kalos}\\
  {\small Asesor: Martinez Abaunza, Victor Eduardo}\\
  {\small Lima, 15 de Abril de 2026}
\end{center}

\vspace{4pt}
\hrule
\vspace{8pt}

%% ---------------------------------------------------------------
\section{Título del Proyecto}

Mitigación de la Resistencia Evolutiva en Glioblastoma Multiforme mediante
MARL Adversarial: Optimización de Políticas Adaptativas bajo un Framework
MAPPO-CTDE y Simulación Calibrada.

%% ---------------------------------------------------------------
\section{Planteamiento del Problema}

Los protocolos de dosificación estáticos, lineales y agresivos son incapaces de adaptarse a la
dinámica evolutiva del cáncer. Al aplicar dosis fijas, la selección natural
elimina las células sensibles y favorece la proliferación de clones resistentes,
reduciendo drásticamente la ventana terapéutica del paciente.

Desde la perspectiva computacional, este es un problema de \textbf{optimización
dinámica en un entorno adversarial}: el sistema (el tumor) modifica sus reglas
de transición de estado en respuesta a cada acción terapéutica. Las estrategias
estáticas convergen a mínimos locales donde la recompensa del tratamiento cae a
cero, equivalente a la resistencia clínica completa.

La propuesta es un \textbf{Framework de Simulación Co-evolutiva} basado en
Aprendizaje por Refuerzo Multi-Agente (MARL) con el esquema MAPPO-CTDE
(Centralized Training, Decentralized Execution). Dos agentes compiten: el
\textit{Agente Terapia} optimiza la dosificación adaptativa, mientras que el
\textit{Agente Cáncer} actúa como adversario que modela la presión evolutiva,
forzando al primero a anticipar la resistencia en lugar de solo reaccionar.

%% ---------------------------------------------------------------
\section{Objetivo General}

Diseñar e implementar un framework de simulación co-evolutiva basado en MARL
adversarial (MAPPO-CTDE) que aprenda políticas de dosificación adaptativa para
el tratamiento de Glioblastoma Multiforme, minimizando la aparición de
resistencia farmacológica mediante la anticipación de la dinámica evolutiva
tumoral, y validando su desempeño frente al protocolo estándar de dosis máxima
tolerada sobre un entorno de simulación calibrado con datos biológicos reales.

%% ---------------------------------------------------------------
\section{Trabajos Relacionados}

\begin{enumerate}[leftmargin=*, itemsep=3pt, topsep=2pt]

  \item \textbf{Kratz \& Adamczyk (2025)} -- \textit{Reinforcement Learning for
  Control of Non-Markovian Cellular Population Dynamics}. ICLR 2025 /
  NeurIPS ML4PS Workshop 2024. arXiv:2410.08439. El trabajo más cercano al
  propuesto: un agente RL entrena en un entorno Gymnasium de poblaciones
  celulares sensibles/resistentes con dinámica no-markoviana, aprendiendo
  protocolos de dosificación adaptativos. Valida la viabilidad del
  enfoque RL en este dominio y sirve como baseline metodológico directo.

  \item \textbf{Mashayekhi et al. (2024)} -- \textit{Deep Reinforcement
  Learning-Based Control of Chemo-Drug Dose in Cancer Treatment}.
  \textit{Comput. Methods Programs Biomed.}, 243, 107884. Controlador DRL
  para dosificación de quimioterapia evaluado sobre pacientes simulados con
  espacios de estado y acción continuos; referencia de comparación directa
  (agente único) que nuestro framework adversarial busca superar.

  \item \textbf{Li et al. (2025)} -- \textit{Optimal Adaptive Cancer Therapy
  Based on Evolutionary Game Theory}. \textit{PLOS ONE}, 20(4): e0320677.
  Incorpora farmacocinética en un modelo de teoría de juegos evolutivos e
  identifica estrategias óptimas de terapia adaptativa mediante el Principio
  de Mínimo de Pontryagin. Fundamenta matemáticamente el modelo de
  competencia celular calibrado con GDSC y ecuaciones de Lotka-Volterra.

  \item \textbf{Salvioli et al. (2024)} -- \textit{Stackelberg Evolutionary
  Games of Cancer Treatment}. \textit{Dynamic Games Appl.}
  DOI:10.1007/s13235-024-00609-z. Modela el tratamiento oncológico como un
  juego de Stackelberg médico--cáncer donde el tumor responde evolucionando
  resistencia; establece la base teórica del esquema adversarial que
  implementamos con MARL en lugar de control óptimo clásico.

  \item \textbf{Yu et al. (2022)} -- \textit{The Surprising Effectiveness of
  PPO in Cooperative Multi-Agent Games (MAPPO)}. NeurIPS 2022.
  arXiv:2103.01955. Paper fundacional del algoritmo MAPPO-CTDE adoptado como
  arquitectura central del framework: demuestra que PPO con crítico
  centralizado logra rendimiento superior o comparable a métodos off-policy
  más complejos en entornos multi-agente, con mínimo ajuste de
  hiperparámetros.

\end{enumerate}

%% ---------------------------------------------------------------
\section{Posibles Datasets}

\begin{center}
\small
\begin{tabular}{>{\bfseries}p{3.2cm} p{4.8cm} p{3.2cm}}
\toprule
Tipo de dato & Fuente & Uso en el proyecto \\
\midrule
Sensibilidad química &
  GDSC / NCI-60 &
  Calibrar el daño del fármaco sobre cada línea celular \\
Perfiles ómicos &
  CCLE (Cancer Cell Line Encyclopedia) &
  Definir el espacio de estados genómicos del tumor \\
Estructura molecular &
  SMILES via PubChem &
  Representar fármacos para el Agente Terapia \\
Dinámica poblacional &
  Modelos ODE (Lotka-Volterra, generado) &
  Simular competencia células sensibles vs.\ resistentes \\
\bottomrule
\end{tabular}
\end{center}

{\footnotesize Todos los datasets listados son de acceso público y se
encuentran disponibles para descarga directa en sus repositorios oficiales.}

\end{document}